\section{工程实现细节}

\subsection{项目结构与模块职责}

当前仓库的设计并不是单纯为了完成一次分类实验，而是围绕“可复现、可扩展、可检查”的目标进行组织。整个项目被划分为数据层、特征层、模型层、可视化层和实验入口层五个部分。这样的模块化结构有两个直接好处：其一，便于不同阶段独立调试；其二，便于在论文中将“算法思想”和“工程实现”分开叙述。

数据层主要负责读取 RNA、甲基化和临床标签，并在加载阶段完成样本对齐和标签规范化。特征层负责对高维输入进行筛选和缩放，并将 MOFA 作为一种跨组学表示学习方法引入管线。模型层仅保留一个统一的 SVM 分类器，目的是把比较重点放在数据融合方式上，而不是分类器差异上。实验入口层则通过 YAML 配置文件驱动不同实验分支，使实验策略从代码中抽离出来，降低修改成本。

\subsection{数据加载与标签清洗}

从工程角度看，数据加载是整个系统最关键也是最容易被忽略的部分。临床标签文件通常并不总是满足统一格式要求，尤其在不同下载源或手工整理版本中，分隔符、列名、重复行和标签命名都可能发生变化。为此，当前实现中在加载阶段显式读取 TSV 文件，并完成列名归一化、标签重命名、重复样本删除和亚型过滤。

这一步看似琐碎，实际上决定了后续所有实验能否站在同一块基石上。如果样本对齐有误，那么后面再复杂的模型都没有意义。与许多“只谈算法不谈数据”的写法不同，本文特别强调这一点，是因为多组学研究的失败常常并非模型失败，而是数据入口失败。

\subsection{样本对齐的实现逻辑}

对齐模块的设计目标是让 RNA、甲基化和临床标签都使用同一组样本，并且保持稳定顺序。当前实现采用基于 RNA 样本列顺序的交集筛选方式，原因是 RNA 矩阵通常是实验的主输入，保留其原始顺序能够让后续结果更容易追踪。

在数学上，这个过程可以写成集合交集 $S = S_r \cap S_m \cap S_c$，其中 $S_r$、$S_m$ 和 $S_c$ 分别代表 RNA、甲基化和临床可用样本集合。随后，我们将最终顺序固定为 RNA 的原始列顺序中属于 $S$ 的那部分样本。这样的设计既满足样本一致性，又避免了集合操作产生的不确定顺序。

\subsection{预处理器的工程化封装}

为了避免泄露并保持特征维度一致性，本文将 RNA 与甲基化预处理器封装成带有 \texttt{fit} 与 \texttt{transform} 接口的类。这里的关键不是“把代码写成类”本身，而是让特征筛选、统计量拟合与数据变换三个环节在工程上分离。

RNA 预处理器包含对数变换、均值过滤、方差筛选和 z-score 标准化。甲基化预处理器则更直接一些，主要是缺失值删除、高方差筛选和标准化。通过把特征列表存储为对象内部状态，后续训练集和测试集就可以共享同一组特征索引，从而保证输入一致。

\subsection{MOFA 接口适配}

MOFA 是整个工程中最容易出问题的环节。其原因并不在于模型思想复杂，而在于第三方库 \texttt{mofapy2} 的输入格式和训练顺序要求较严格。当前实现已经解决了三个关键问题：一是把数据整理为视图×分组的嵌套列表；二是在 \texttt{build()} 前显式调用 \texttt{set\_train\_options()}；三是在输入维度较小时避免训练阶段的初始化失败。

更重要的是，当前实现将 MOFA 的输出当作下游分类器的输入，而不是直接作为最终分类器。这一点在论文里非常值得强调，因为它把“多组学表示学习”与“分类判别”两个任务分开处理，使方法链条更加清晰。

\subsection{实验脚本与配置文件}

实验脚本的意义在于把研究结论从“可复现”推进到“可一键运行”。主实验脚本 \texttt{run\_all.sh} 用于跑 RNA、Concat 和 MOFA 三组对照；消融脚本 \texttt{scripts/run\_ablation.sh} 用于集中验证去甲基化、去 RNA 与不做特征选择三类设置。对于论文来说，这类脚本说明研究不是一次性手工实验，而是可重复执行的实验流程。

YAML 配置文件进一步把实验参数显式化。这样做的好处是，在论文正文中可以把“实验设置”写成表格和文字说明，而在代码层面则通过配置文件保证参数一致性。这种“配置驱动”的工程方式，特别适合课程设计和研究报告，因为它使实验过程更容易审阅和复查。
